\section{Data, Preprocessing, and Reproducibility Details}
\label{si:legacy_details}

This section preserves and expands the data and implementation information from
the previously prepared Supporting Information. It is included so that the new
equation guide and standalone artwork do not remove any reproducibility details
from the earlier SI version.

\subsection{Dataset Sources and Processed Splits}

We evaluate four public drug-target interaction benchmarks: BindingDB,
BioSNAP, Human, and C. elegans. Each sample contains a drug SMILES string, a
protein sequence, and a binary interaction label. The statistics below are
computed from the processed training, validation, and test split files used by
the experiments. The split files are fixed across repeated runs; the seed is
varied for model initialization and training stochasticity.

\begin{table}[!htbp]
\centering
\caption{Processed dataset statistics used by the manuscript.}
\label{tab:si_processed_dataset_stats}
\begin{tabular}{lrrrrrr}
\toprule
Dataset & Drugs & Proteins & Interactions & Positive & Negative & Positive rate \\
\midrule
BindingDB & 14,643 & 2,623 & 49,199 & 20,674 & 28,525 & 42.02\% \\
BioSNAP & 4,505 & 2,181 & 27,464 & 13,830 & 13,634 & 50.36\% \\
Human & 2,726 & 2,001 & 5,997 & 2,633 & 3,364 & 43.91\% \\
C. elegans & 1,767 & 1,876 & 7,786 & 3,893 & 3,893 & 50.00\% \\
\bottomrule
\end{tabular}
\end{table}

\begin{table}[!htbp]
\centering
\caption{Training-split class counts and positive-class weights.}
\label{tab:si_training_weights}
\begin{tabular}{lrrr}
\toprule
Dataset & Training positives & Training negatives & Positive weight $N_-/N_+$ \\
\midrule
BindingDB & 14,583 & 19,856 & 1.3616 \\
BioSNAP & 9,684 & 9,540 & 0.9851 \\
Human & 1,846 & 2,351 & 1.2736 \\
C. elegans & 2,727 & 2,723 & 0.9985 \\
\bottomrule
\end{tabular}
\end{table}

Data access points are listed in the manuscript Data Availability Statement:
BindingDB, BioSNAP, and the Human/C. elegans benchmark files in the
TransformerCPI repository.

\subsection{Input Processing and Feature-Cache Schema}

SMILES strings are encoded by a ChemBERTa hidden-state cache and converted to
DGL molecular graphs. Protein sequences are encoded by a ProtBERT hidden-state
cache. The data loader validates cache shapes and retains explicit masks for
protein residues and invalid conformer slots. The main model uses the vanilla
cache variant, \texttt{drug3d\_features.pt}.

\begin{table}[!htbp]
\centering
\caption{Feature-cache tensors and dimensions.}
\label{tab:si_cache_schema}
\scriptsize
\setlength{\tabcolsep}{3pt}
\begin{tabular}{>{\RaggedRight\arraybackslash}p{0.15\linewidth}>{\RaggedRight\arraybackslash}p{0.24\linewidth}>{\RaggedRight\arraybackslash}p{0.18\linewidth}>{\RaggedRight\arraybackslash}p{0.32\linewidth}}
\toprule
Stream & File or tensor & Shape & Processing and use \\
\midrule
Drug 1D & \texttt{smiles\_features.pt} & $354\times128$ & ChemBERTa token hidden states after adaptive pooling. \\
Drug graph & DGL graph & Up to $290\times75$ & Canonical atom/bond features; self-loops enabled. \\
Drug 3D & \texttt{drug3d\_features.pt} & $K\times64\times128$ & Per-conformer atom features; $K=8$ in the main model. \\
Coordinates & 3D conformer coordinates & $K\times64\times3$ & Coordinates consumed by the EGNN branch. \\
Conformer metadata & \texttt{conf\_mask}, energy & $K$, $K$ & Validity mask and UFF energy for target-aware scoring. \\
Protein & \texttt{protein\_features.pt} & $128\times1024$ & ProtBERT residue features after pooling; residue mask retained. \\
\bottomrule
\end{tabular}
\end{table}

RDKit ETKDGv3 generates up to eight conformers. The extraction routine uses
random seed 42, RMSD pruning threshold 0.5, and one embedding thread. UFF
optimization is attempted for at most 100 iterations when parameters are
available. Molecules with more than 160 heavy atoms are skipped by the 3D
extraction guard. Protein sequences are upper-cased after whitespace removal;
U, Z, O, and B are replaced with X before token encoding. The model-side
protein length is padded or masked to 128. The feature-extraction schema
version is 3.

\subsection{Complete Architecture Configuration}

The model uses a 128-dimensional shared latent space. Protein context first
conditions target-aware conformer scoring. The resulting drug representation
provides a ligand context for FiLM modulation of protein tokens. A
bidirectional protein-side Mamba block refines the protein sequence before
one-layer BiIntention cross-modal matching.

\begin{table}[!htbp]
\centering
\caption{Main TAMR-DTI architecture configuration.}
\label{tab:si_architecture_config}
\scriptsize
\setlength{\tabcolsep}{4pt}
\begin{tabular}{>{\RaggedRight\arraybackslash}p{0.25\linewidth}>{\RaggedRight\arraybackslash}p{0.64\linewidth}}
\toprule
Component & Setting \\
\midrule
Drug graph encoder & Input node dimension 75; hidden dimensions [128, 128, 128]; maximum graph length 290. \\
Target-aware conformer encoder & EGNN dimension 128; $K=8$; 64 atoms per conformer; protein-conditioned score input; energy projection included. \\
Drug fusion & Graph and conformer tokens form 354 structural slots; layer-normalized fusion with a learned sigmoid gate initialized at bias $-2.0$. \\
Protein encoder & 1024-dimensional ProtBERT features projected to 128; three ProteinACmix layers; kernels 3, 6, 9; eight heads; ligand FiLM scale 0.05. \\
Protein Mamba & Protein tokens only; bidirectional; $d_{\rm state}=16$; $d_{\rm conv}=4$; expansion 2; dropout 0.1; Mamba-1 backend; gate bias $-2.0$. \\
BiIntention & One layer; eight heads; embedding dimension 128; bidirectional drug-protein intention updates. \\
Prediction head & Four-layer MLP; input 256; hidden 256; intermediate output 128; binary logit output. \\
\bottomrule
\end{tabular}
\end{table}

\subsection{Training, Loss, and Evaluation}

\begin{table}[!htbp]
\centering
\caption{Training and evaluation settings.}
\label{tab:si_training_evaluation}
\scriptsize
\setlength{\tabcolsep}{4pt}
\begin{tabular}{>{\RaggedRight\arraybackslash}p{0.28\linewidth}>{\RaggedRight\arraybackslash}p{0.59\linewidth}}
\toprule
Setting & Value \\
\midrule
Optimizer & Adam \\
Initial learning rate & $1\times10^{-4}$ \\
Weight decay & $1\times10^{-4}$ \\
Maximum epochs & 100 \\
Batch and accumulation & 16 per GPU; gradient accumulation 4 steps \\
Gradient clipping & Maximum norm 1.0 \\
Scheduler & Plateau; factor 0.5; patience 5; minimum learning rate $1\times10^{-6}$ \\
Early stopping & Patience 10; minimum improvement $1\times10^{-4}$ \\
Checkpoint & Highest validation AUROC \\
Runtime & Python 3.9; PyTorch, DeepSpeed ZeRO-2, DGL/DGLLife, RDKit, Transformers, and mamba-ssm; eight Tesla P100 GPUs. \\
\bottomrule
\end{tabular}
\end{table}

For completeness, the threshold-dependent metrics used in the manuscript are
defined below. The threshold $\tau$ is applied to the predicted probability.
\begin{align}
\mathrm{Accuracy}(\tau)
 &= \frac{\mathrm{TP}(\tau)+\mathrm{TN}(\tau)}
 {\mathrm{TP}(\tau)+\mathrm{TN}(\tau)+\mathrm{FP}(\tau)+\mathrm{FN}(\tau)},\\
\mathrm{Precision}(\tau)
 &= \frac{\mathrm{TP}(\tau)}{\mathrm{TP}(\tau)+\mathrm{FP}(\tau)},\\
\mathrm{Recall}(\tau)
 &= \frac{\mathrm{TP}(\tau)}{\mathrm{TP}(\tau)+\mathrm{FN}(\tau)},\\
\mathrm{F1}(\tau)
 &= \frac{2\,\mathrm{Precision}(\tau)\,\mathrm{Recall}(\tau)}
 {\mathrm{Precision}(\tau)+\mathrm{Recall}(\tau)}.
\end{align}
The validation threshold is searched over $0.01,0.02,\ldots,0.99$ and then
applied unchanged to the held-out test split. Fixed-threshold results use
$\tau=0.5$.

\subsection{Main Run Matrix}

All four datasets use the same main architecture and five fixed seeds, 42
through 46. Configuration identifiers correspond to files under
\texttt{configs/experiments/}.

\begin{table}[!htbp]
\centering
\caption{Main repeated-seed configuration index.}
\label{tab:si_run_matrix}
\scriptsize
\setlength{\tabcolsep}{5pt}
\begin{tabular}{lll}
\toprule
Dataset & Seed & Configuration \\
\midrule
BindingDB & 42 & \texttt{stage2\_main\_01} \\
BindingDB & 43 & \texttt{stage2\_main\_02} \\
BindingDB & 44 & \texttt{stage2\_main\_03} \\
BindingDB & 45 & \texttt{stage2\_main\_04} \\
BindingDB & 46 & \texttt{stage2\_main\_05} \\
BioSNAP & 42 & \texttt{stage2\_main\_06} \\
BioSNAP & 43 & \texttt{stage2\_main\_07} \\
BioSNAP & 44 & \texttt{stage2\_main\_08} \\
BioSNAP & 45 & \texttt{stage2\_main\_09} \\
BioSNAP & 46 & \texttt{stage2\_main\_10} \\
Human & 42 & \texttt{stage2\_main\_11} \\
Human & 43 & \texttt{stage2\_main\_12} \\
Human & 44 & \texttt{stage2\_main\_13} \\
Human & 45 & \texttt{stage2\_main\_14} \\
Human & 46 & \texttt{stage2\_main\_15} \\
C. elegans & 42 & \texttt{stage2\_main\_16} \\
C. elegans & 43 & \texttt{stage2\_main\_17} \\
C. elegans & 44 & \texttt{stage2\_main\_18} \\
C. elegans & 45 & \texttt{stage2\_main\_19} \\
C. elegans & 46 & \texttt{stage2\_main\_20} \\
\bottomrule
\end{tabular}
\end{table}

\subsection{Expanded Performance and Diagnostic Results}

The standalone tables in Section~\ref{si:tables} reproduce the main
AUROC/AUPRC comparison, the repeated-run summary, the BioSNAP module ablation,
and the conformer-count diagnostic. All diagnostics use BioSNAP, seed 42, the
same fixed split, validation-AUROC checkpoint selection, and fixed-threshold
evaluation. They are single-seed mechanism diagnostics rather than repeated-run
significance tests.

\subsection{Conformer-Weight and Mamba Interpretability}

The analysis uses all 5,493 BioSNAP test pairs from the seed-42 target-aware
model. For each sample, conformer weights are sorted in descending order, and
prediction groups use validation threshold 0.62. Normalized conformer entropy is
calculated as
\begin{equation}
\bar H_i=-\frac{1}{\log K_i}\sum_{k=1}^{K_i}w_{ik}\log(w_{ik}+\epsilon),
\end{equation}
where $K_i$ is the number of valid conformers for sample $i$ and $\epsilon$ is
a numerical-stability constant.

\begin{table}[!htbp]
\centering
\caption{Population-level interpretability statistics.}
\label{tab:si_interpretability_statistics}
\scriptsize
\setlength{\tabcolsep}{4pt}
\begin{tabular}{lrrrr}
\toprule
Group & $n$ & Top-1 weight & Entropy & Mamba refine mean \\
\midrule
All test pairs & 5,493 & 0.5975 & 0.4799 & 0.094526 \\
True positive & 2,175 & $0.5910\pm0.2341$ & $0.4876\pm0.2545$ & 0.093963 \\
True negative & 2,328 & $0.5974\pm0.2311$ & $0.4808\pm0.2517$ & 0.095077 \\
False positive & 417 & $0.5857\pm0.2287$ & $0.4834\pm0.2468$ & 0.093675 \\
False negative & 573 & $0.6309\pm0.2401$ & $0.4446\pm0.2651$ & 0.095047 \\
\bottomrule
\end{tabular}
\end{table}

\begin{table}[!htbp]
\centering
\caption{Median sorted conformer weights by prediction group.}
\label{tab:si_sorted_conformer_weights}
\begin{tabular}{lrrrr}
\toprule
Group & Rank 1 & Rank 2 & Rank 3 & Rank 4 \\
\midrule
True positive & 0.5320 & 0.2479 & 0.0946 & 0.0121 \\
True negative & 0.5288 & 0.2457 & 0.0964 & 0.0137 \\
False positive & 0.5192 & 0.2579 & 0.1036 & 0.0132 \\
False negative & 0.5804 & 0.2317 & 0.0756 & 0.0053 \\
\bottomrule
\end{tabular}
\end{table}

\begin{table}[!htbp]
\centering
\caption{Auxiliary modulation statistics from the interpretability analysis.}
\label{tab:si_auxiliary_modulation}
\begin{tabular}{lr}
\toprule
Quantity & Mean \\
\midrule
Mamba gate & 0.132949 \\
Conformer top-1 margin & 0.365227 \\
Protein FiLM layer 1 strength & 0.020351 \\
Protein FiLM layer 2 strength & 0.013285 \\
Protein FiLM layer 3 strength & 0.012323 \\
\bottomrule
\end{tabular}
\end{table}

The sample-level analysis file is
\url{outputs/interpretability/biosnap_seed42/biosnap_seed42_interpretability_samples.csv}.
It records prediction groups, threshold, conformer entropy, top-1 weight,
conformer margin, Mamba summaries, eight conformer weights, and layer-wise FiLM
statistics.

\subsection{Reproducibility and Availability}

The implementation, preprocessing code, and experiment configuration files are
available at \url{https://github.com/leleleocc/TAMR-DTI}. Model code is under
\texttt{src/}, feature extraction is under \texttt{scripts/pre\_extract.py},
and the run configurations listed in Table~\ref{tab:si_run_matrix} are under
\texttt{configs/experiments/}. To reproduce a run, obtain the public benchmark
files, construct the fixed split directory with \texttt{train.csv},
\texttt{val.csv}, and \texttt{test.csv}, build the ChemBERTa, ProtBERT, graph,
and 3D caches with the settings above, and launch the corresponding
configuration. Select the checkpoint by validation AUROC and evaluate the
held-out test set under the metric conventions in this SI and the manuscript.
